\chapter{Semantic, Memory, and Effect-Aware IRs}
\label{ch:semantic-effect-ir}

Two adjacent stores look independent when they use different variable names.
They cease to be independent when those variables hold the same address.
Two additions look interchangeable when they compute the same mathematical integer.
They cease to be interchangeable when one wraps and the other traps.
An IR must expose enough meaning for the optimizer to tell these cases apart.

An opcode list alone is therefore not a semantic specification.
The contract must also define values, state, observable behavior, exceptional outcomes, and the assumptions under which a transformation may replace an operation.
This chapter connects that contract to explicit memory dependencies and effect-aware scheduling.

\section{Define an operation as a state transition}
\label{sec:ch63-semantics}

For a small sequential language, let a state contain local values, a heap, and an observable-event sequence.
Evaluating an operation either produces a new state and result, raises a specified error, or fails to terminate.
The event sequence might record output, foreign calls, or accesses designated observable by the language.
Ordinary internal loads need not appear in it; their returned values and possible faults still matter.

A \term{semantic IR} gives these transitions a precise interpretation rather than relying on the implementation of the current interpreter as an undocumented definition.
Executable semantics can power interpreters, tests, validators, and instruction selection, but each generated tool must implement that interpretation correctly.
One specification does not by itself prove that all its consumers agree.

Consider a bounds-checked array read followed by output.
Moving output before the read changes the trace if the read fails.
Removing the read changes termination if it can block forever.
Even when its numeric result is dead, the operation may carry required behavior.
Table~\ref{tab:ch63-effects} distinguishes the relevant dimensions.

\begin{table}[htbp]
\caption{Effect summaries describe behavior beyond the returned value.}
\label{tab:ch63-effects}
\centering
\begin{tabular}{ll}
\toprule
Component & Meaning in the sequential model \\
\midrule
Read set & memory regions whose contents may be inspected \\
Write set & memory regions whose contents may change \\
Observable & an externally visible event may occur \\
May trap & an exceptional outcome may terminate the operation \\
May diverge & the operation may not return \\
\bottomrule
\end{tabular}
\end{table}

\section{Make memory versions explicit}
\label{sec:ch63-memory}

One simple IR threads a memory-state token through every store and effectful call.
A load consumes the memory version it observes; a store produces a new version.
Control-flow joins merge incoming versions with a memory merge node.
This exposes sequencing to graph algorithms without pretending that a memory token is a runtime pointer.
Figure~\ref{fig:ch63-token} shows a store to $p$, a store to $q$, and a load from $p$.

\begin{figure}[htbp]
\centering
\begin{tikzpicture}[node distance=13mm]
\node[basic block] (m0) {memory m0};
\node[basic block,below=of m0] (s1) {m1 = store(m0,p,9)};
\node[basic block,below=of s1] (s2) {m2 = store(m1,q,4)};
\node[basic block,below=of s2] (load) {v = load(m2,p)};
\draw[cfg edge] (m0)--(s1); \draw[cfg edge] (s1)--(s2); \draw[cfg edge] (s2)--(load);
\end{tikzpicture}
\caption{Memory-state edges preserve order until alias and effect reasoning justify bypassing a dependency.}
\label{fig:ch63-token}
\end{figure}

A single chain conservatively orders operations that may actually be independent.
If $p$ and $q$ denote distinct ordinary cells, the second store cannot change the later load from $p$.
The load can then be replaced by $9$, subject to the access-validity contract.
If $p=q$, its value is $4$ instead.
Table~\ref{tab:ch63-alias} works both cases, including the result of swapping the stores.

\begin{table}[htbp]
\caption{Alias information changes both load forwarding and reorder legality.}
\label{tab:ch63-alias}
\centering
\begin{tabular}{llll}
\toprule
Addresses & Original load & Swapped load & Swap legal here \\
\midrule
$p\ne q$ & $9$ & $9$ & yes \\
$p=q$ & $4$ & $9$ & no \\
\bottomrule
\end{tabular}
\end{table}

LLVM's MemorySSA is an analysis overlay describing memory definitions, uses, and merges; it is not a complete executable heap semantics or a claim that alias analysis has become unnecessary~\cite{llvm-2026-memoryssa}.
Its walker uses alias information to find relevant clobbers.
Firm provides a graph-based design in which memory dependencies are part of the IR's organization~\cite{braun-2011-firm}.
These designs share a motivation while making different engineering choices about where dependency information lives and how it is updated.

\section{Use effect summaries as conservative permissions}
\label{sec:ch63-commute}

For an operation $a$, let $R_a$ and $W_a$ overapproximate its read and write regions.
A sufficient memory-independence condition for two operations is
\[
W_a\cap(R_b\cup W_b)=\emptyset,
\qquad W_b\cap R_a=\emptyset.
\]
The regions must describe overlapping storage conservatively; distinct syntactic addresses are not enough.
The operations must also have no SSA value dependence in either direction, operate on valid storage, terminate, and produce no required ordered events or faults under this rule.
These are sufficient conditions for our sequential model, not a complete concurrency memory model.

Algorithm~\ref{alg:ch63-commute} is deliberately conservative.
For hashed explicit region sets, intersections can be checked in expected time linear in the input set sizes.
A region analysis may instead use intervals, allocation identities, or a top element meaning unknown memory.
Treating unknown as an empty set would turn uncertainty into an unsound permission.

\begin{algorithm}[htbp]
\caption{Check a sufficient condition for swapping two sequential operations.}
\label{alg:ch63-commute}
\begin{pseudocode}
\Require sound effect summaries for $a,b$, valid accesses, no value dependence
\Ensure true only when this rule permits exchanging their order
\If{either operation may trap, diverge, or emit an observable event}
  \State \Return false
\EndIf
\If{$W_a\cap(R_b\cup W_b)\ne\emptyset$ or $W_b\cap R_a\ne\emptyset$}
  \State \Return false
\EndIf
\State \Return true
\end{pseudocode}
\end{algorithm}

The companion tests every pair of four concrete addresses against execution of the store example.
This checks the concrete disjointness instance, not a compiler's entire alias analysis.
Exceptions, atomics, volatile accesses, synchronization, and signals require richer contracts before this rule can be generalized.
In particular, absence of ordinary read/write overlap does not authorize moving a synchronization operation.

\section{Effects and types serve different roles}
\label{sec:ch63-types}

A type may say that an operation returns an integer without saying whether it writes a global variable.
An effect annotation adds that information, either as part of the type or in a separate checked summary.
A \term{monadic IR} makes sequencing explicit through its computation structure, allowing value construction and effectful computation to be distinguished.
MIL demonstrates such an organization for strict functional languages and connects block interfaces to specialization and representation changes~\cite{jones-2018-mil}.

The implementation must still decide how summaries are obtained.
Trusted declarations are a language boundary; inferred summaries need a sound analysis; unchecked annotations are merely hints.
Function summaries for recursion typically require a fixed point.
A caller that specializes against an imported effect summary must invalidate its result if that summary changes.
This is one concrete dependency for the query architecture of Chapter~\ref{ch:incremental-databases}.
Foreign callbacks make this dependency transitive: a native visitor that performs no direct stores can still modify managed globals through its callback.
Chapter~\ref{ch:ffi-foundations} develops a boundary-admission contract that includes callbacks, pointer retention, and failure paths before an imported summary can justify optimization.

\section{Carry contracts across lowering}
\label{sec:ch63-lowering}

Lowering a bounds-checked load to address arithmetic plus a machine load introduces obligations.
The index calculation must use the intended integer width; the check must dominate the access; pointer arithmetic must stay within the required object model; and the failure branch must retain the specified error.
A semantic operation may become several instructions, but its observable contract remains.
Figure~\ref{fig:ch63-lowering} separates the value, memory, and exceptional components.

\begin{figure}[htbp]
\centering
\begin{tikzpicture}[node distance=14mm]
\node[basic block] (source) {checked array read};
\node[basic block,below=of source] (guard) {bounds test};
\node[basic block,below left=of guard] (error) {specified error};
\node[basic block,below right=of guard] (load) {address + load};
\draw[cfg edge] (source)--(guard);
\draw[cfg edge] (guard)--node[left,font=\small] {fail}(error);
\draw[cfg edge] (guard)--node[right,font=\small] {pass}(load);
\end{tikzpicture}
\caption{Lowering preserves exceptional behavior as well as the successful value.}
\label{fig:ch63-lowering}
\end{figure}

Algorithm~\ref{alg:ch63-contract} describes a validation boundary, not a decision procedure for arbitrary program equivalence.
Local proof rules, a solver in a decidable fragment, or finite tests have different evidential strength.
An unsupported obligation must remain unresolved rather than being interpreted as a successful proof.

\begin{algorithm}[htbp]
\caption{Validate a proposed semantic lowering through explicit obligations.}
\label{alg:ch63-contract}
\begin{pseudocode}
\Require source operation contract, target sequence, established input facts
\Ensure accepted lowering only when every required obligation is discharged
\State check target typing, widths, control edges, and memory dependencies
\State formulate successful-value and final-memory obligations
\State formulate required event, exceptional-outcome, and termination obligations
\ForAll{required obligations}
  \If{the selected proof mechanism cannot establish the obligation}
    \State \Return unresolved or rejected
  \EndIf
\EndFor
\State \Return accepted with recorded assumptions
\end{pseudocode}
\end{algorithm}

There is no general finite complexity bound for discharging these semantic obligations: unrestricted equivalence is undecidable.
The structural checks are finite graph checks, while the proof mechanism determines the remaining cost and completeness.
Keeping that boundary visible prevents a useful executable specification from being mistaken for an automatic proof of every optimization.

\section{Exercises}
\label{sec:ch63-exercises}

\begin{enumerate}
\item\label{ex:ch63-stores} Execute the two stores and load for $p=1,q=2$, and then for $p=q=1$.
Give both final memory maps and explain the swap decision.
\item\label{ex:ch63-effects} Operation $a$ reads region $A$ and writes $B$; operation $b$ reads $B$ and writes $C$.
Assume pairwise disjoint regions and no other effects.
Does the sufficient swap rule accept? Identify the dependency.
\item\label{ex:ch63-trap} Construct a program where a dead result cannot be removed because computing it can trap before an output operation.
State what additional fact would permit removal.
\item\label{ex:ch63-summary} A callee's summary changes from reading $A$ to writing $A$.
Explain which cached caller optimization must be invalidated and why a type-only dependency is insufficient.
\end{enumerate}

\section{Further reading}
\label{sec:ch63-reading}

LLVM's documentation explains MemorySSA's access nodes and clobber queries~\cite{llvm-2026-memoryssa}.
Firm and MIL offer contrasting graph and monadic approaches to exposing dependencies and representation choices~\cite{braun-2011-firm,jones-2018-mil}.
The executable \code{Effect} model in \path{example-code/architecture-labs/labs.py} isolates the small sequential commutation rule used here.
